\documentclass[11pt]{article}
\usepackage[margin=0.85in]{geometry}
\usepackage[english]{babel}
\usepackage{amsmath,amssymb,array,booktabs,tabularx}
\usepackage[hidelinks]{hyperref}

\setlength{\parindent}{0pt}
\setlength{\parskip}{0.55\baselineskip}
\setlength{\emergencystretch}{3em}
\newcolumntype{L}[1]{>{\raggedright\arraybackslash}p{#1}}

\title{Technical Review of\\
\emph{Public Emergency Codes: Evidence Audit, Assumption-Conditional Mortality
Scenarios, and Full-Implementation Potential}}
\author{}
\date{August 8, 2026}

\begin{document}
\maketitle

\section*{Overall assessment}

The manuscript is unusually explicit that its PEC mortality outputs are
assumption-conditional planning quantities rather than identified causal
effects.  The displayed modal calculations were checked against the supplied
JSON files: the current-attribution total ($67.416$), mature all-pathway total
($2657.158$), and Tier-C-zero total ($1352.812$) agree before rounding.  Units
and benefit-minus-harm signs are also internally consistent in the displayed
equations.

Major revision is nevertheless required before the mortality headlines are
technically defensible.  The current OHCA calculation combines a whole-registry
denominator and survival probability with effect slopes from narrower clinical
populations; the mature model does not mathematically enforce its claimed
exclusion between the separate cardiac-arrest row and the acute language/access
rows.  The Tier-C-zero result also remains dominated by positive PEC-specific
quantities that the manuscript itself identifies as author-specified priors.
The supplied code reproduces the archived results exactly, but its
version-control/deposit status and machine-readable seed metadata need repair.
These are structural issues, not presentation details.

\section*{Technical findings}

\subsection*{1. An unidentified effect is assigned the numerical value zero}

\textbf{Status: definite conceptual error.}

\textbf{Location.} Abstract; ``Scope, estimands, and interpretation'';
``Observed evidence without PEC attribution''; and the boxed statement
\[
\boxed{L_{\mathrm{ID}}=0\text{ identified lives/year}.}
\]

\textbf{Problem.} The paper first defines $L_{\mathrm{ID}}$ as a PEC mortality
effect identified by current data, then states that $L_{\mathrm{ID}}$ is not
identified, and finally assigns it the value zero.  Failure to identify or
estimate an effect is not an estimate of a zero effect.  The qualifier
``identified lives'' explains the intended bookkeeping convention but does not
repair the estimand.

\textbf{Why it matters.} The notation can be read, cited, or imported into a
downstream model as an observed null.  It also conflicts with the conclusion's
correct statement that the effect ``remains unidentified.''

\textbf{Suggested fix.} Report the evidence-only PEC effect as ``not identified''
or ``not estimable from current data,'' without a numerical equality.  If a
bookkeeping quantity is needed, define a separate variable, for example
\[
L_{\mathrm{credited}}^{\mathrm{evidence\ only}}=0,
\]
and state that this is a decision rule assigning no mortality credit, not an
estimate of the true effect.

\subsection*{2. Equation~\textup{(7)} mixes incompatible OHCA populations}

\textbf{Status: definite model/estimand mismatch; direction and magnitude of
the resulting bias are not identifiable from the manuscript.}

\textbf{Location.} ``Adult OHCA time-response conversion,''
``Incident-level language/access overlap,'' Table~4, and
\begin{align*}
L_{\mathrm{OHCA}}=N_{\mathrm{OHCA}}\Bigl\{&(a_L-a_J)
d(s,\beta_{\mathrm{CPR}},t_L)+(a_A-a_J)d(s,\beta_{\mathrm R},t_A)\\
&+a_J[\operatorname{logit}^{-1}(\operatorname{logit}s
+\beta_{\mathrm{CPR}}t_L+\beta_{\mathrm R}t_A)-s]\Bigr\}.
\end{align*}

\textbf{Problem.} $N_{\mathrm{OHCA}}$ is the full annual EMS-treated
nontraumatic OHCA population and $s=0.105$ is the all-case CARES survival
probability.  In contrast, $\beta_{\mathrm{CPR}}$ is derived from
\texttt{nguyen2024}, whose analyzed population is witnessed OHCA receiving
bystander CPR.  The model has no explicit factor for witnessed status,
bystander/dispatcher CPR eligibility, or receipt of CPR whose start time can be
changed.  The manuscript's cited 2024 CARES summary itself shows that only
$37.5\%$ of registered arrests were bystander-witnessed and that CPR was
initiated by a bystander in $41.9\%$; the relevant intersection is narrower
than either marginal percentage.  Moreover, the $10.5\%$ all-case baseline is
not the baseline survival probability for the slope-estimation subgroup.

\textbf{Why it matters.} Applying a conditional odds-ratio slope to the entire
OHCA denominator is not a valid standardization.  The reported language-only
and joint language/CPR contributions therefore do not correspond to a clearly
defined target population, even though their arithmetic is correct given the
inputs.

\textbf{Suggested fix.} Define a clinically eligible denominator explicitly,
for example
\[
N_{L,\mathrm{eligible}}
=N_{\mathrm{OHCA}}f_{\mathrm{adult}}f_{\mathrm{witnessed}}
f_{\mathrm{CPR\ actionable}},
\]
estimate or elicit each factor transparently, and use a baseline survival
probability from the same target subgroup as the transported slope.  An
alternative is direct standardization over witnessed status, initial rhythm,
location, and CPR source.  Do not hide clinical eligibility inside $u_L$:
Table~4 currently describes $u_L$ as the PEC end-to-end technical/workflow
chain, not as the outcome-study eligibility fraction.

\subsection*{3. The responder-access pathway changes time endpoints without a
transport argument}

\textbf{Status: likely technical issue and required notation/terminology-drift
finding.}

\textbf{Location.} The measured evidence is described as ``building access
time''; $t_A$ is later called ``patient-contact time recovered''; and
$\beta_{\mathrm R}$ is calibrated from ``EMS response-time'' studies in
\texttt{holmen2020} and \texttt{damdin2025}.

\textbf{Problem.} The local \texttt{vertical2007} study measures the interval
from arrival on scene to reaching the patient's side.  The survival studies
used for $\beta_{\mathrm R}$ analyze ambulance/EMS response-time measures in
different systems and populations.  The paper treats these descriptors as if
they identified the same time variable, but it does not specify each clock's
start and stop events or justify transporting the survival association across
those operational phases.

\textbf{Why it matters.} A minute removed before vehicle arrival, a minute
removed while entering a building, and a minute removed before first clinical
intervention need not share the same confounding structure or marginal survival
association.  The zero lower endpoint for $\beta_{\mathrm R}$ represents
uncertainty but does not define the estimand.

\textbf{Suggested fix.} Define separate variables such as
$t_{\mathrm{dispatch\to scene}}$, $t_{\mathrm{scene\to patient}}$, and
$t_{\mathrm{patient\to treatment}}$.  Use an outcome model calibrated to the
same endpoint as the PEC mechanism, or present the current conversion as a
clearly labeled cross-endpoint transport sensitivity with an explicit
transportability assumption.

\subsection*{4. Mature-model exclusion of overlapping OHCA mechanisms is not
mathematically enforced}

\textbf{Status: definite implementation ambiguity and likely double-counting
risk.}

\textbf{Location.} ``Mature full-implementation model,'' Table~8, and
\[
A_j=I_jr_ju_j,\qquad
L_F=\sum_j(\lambda_j^+-\lambda_j^-).
\]

\textbf{Problem.} Four rows based on the $36.367261$ million acute-EMS
denominator are declared mutually exclusive, but ``OHCA bystander assistance''
is outside that group and uses its own $220{,}000/250{,}000/300{,}000$
denominator.  The text says that the OHCA row ``excludes language and
responder-access mechanisms counted elsewhere,'' yet neither the displayed
equations nor the supplied implementation performs this exclusion.  In
\path{pec_full_potential_simulation.py}, normalization is applied only to the
four names in \texttt{acute\_ems\_group}; the cardiac-arrest bystander row is
not a member of that set and is then added separately.  The relevance input
could have been elicited for a pre-excluded subset, but neither the input
definition nor code demonstrates that conditioning.  Nothing in the executable
allocation prevents the same underlying class of cardiac-arrest cases from
being represented in both the separate OHCA row and a language or
location/access share of the acute-EMS denominator.

\textbf{Why it matters.} The manuscript identifies incident-level overlap as a
central correction in the current scenario, but the mature scenario relies on
a verbal exclusion that cannot be checked.  The location/access pathway is the
largest surviving contributor in the Tier-C-zero simulation, so even modest
overlap can materially alter the headline.

\textbf{Suggested fix.} Construct one incident partition that includes the OHCA
submechanisms, or condition the OHCA relevance distribution on not being
assigned to the language/access acute categories.  Publish the precise
allocation formula and a simulation assertion showing that each incident's
mortality outcome is allocated once across all relevant rows, not merely once
within the four-row acute group.

\subsection*{5. ``Tier-C-zero'' is not an evidence-grounded lower scenario}

\textbf{Status: important interpretation problem.}

\textbf{Location.} Abstract, Tables~8--10, ``Results that may be quoted,'' and
the conclusion.

\textbf{Problem.} Zeroing rows labeled C leaves positive effects for
988/crisis, OHCA bystander assistance, silent/language, and location/access.
For every one of these rows, Table~8 and the JSON input file state that the
PEC-specific relevance, implementation, attribution, or absolute risk change
remains author-specified.  Thus the $1746$ median does not show what current
evidence alone supports and should not be interpreted as a conservative
empirical floor.  In addition, every surviving beneficial pathway has
continuous, strictly positive support for its relevance, success, beneficial
ARR, and shared clinical-effect factor.  Near-certain benefit is therefore
largely imposed by the scenario support rather than learned from PEC outcomes.

\textbf{Why it matters.} Readers may infer that the sensitivity removes all
unsupported mortality credit, when it removes only the rows assigned the
letter C.  The label is especially consequential because this result is
prominent in the abstract and companion economic paper.

\textbf{Suggested fix.} Keep the existing sensitivity but label it precisely,
for example ``designated-Tier-C-positive-effects-zero.''  Add a second
``all nonidentified PEC positive effects zero'' analysis that also removes the
positive PEC-specific portions of the external+prior and B+prior rows.  Report
the external association, transported process evidence, PEC attribution, and
PEC ARR priors as separate layers.

\subsection*{6. The numerical results reproduce, but the reproducibility package
and metadata remain incomplete}

\textbf{Status: definite deposit/versioning problem and reporting ambiguity.}

\textbf{Location.} ``Reproducibility,'' ``Complete mature input
distributions,'' and ``Mature probabilistic results.''

\textbf{Problem.} The supplied Python functions were run in memory with the
published draw counts and streams.  Every current and mature archived summary
reproduced exactly (maximum absolute numerical difference $0$), including the
dependence sensitivities and predictive distributions.  However, the
repository's \path{.gitignore} excludes all Python files and the mortality
scripts are not tracked by Git, so an ordinary repository clone or commit can
omit the very code needed for reproduction.  The two load-bearing source
documents cited as \texttt{pecinputs} and \texttt{pecfunctions} are also absent.
Finally, \path{pec_full_potential_results.json} labels \texttt{20260808} simply
as ``seed,'' while the primary RNG is actually initialized with
\texttt{20260809}; the manuscript explains the $+1$ offset, but the
machine-readable metadata does not.  The code also multiplies beneficial ARR
by the shared clinical-effect factor, an operation not shown explicitly in
Equations~(10)--(11).

\textbf{Why it matters.} The current working directory is numerically
reproducible, but the reproducibility claim may fail when the paper is archived
or cloned.  Ambiguous seed metadata also makes it easy for an independent
implementation to use the base seed rather than the actual primary stream.

\textbf{Suggested fix.} Track or otherwise deposit the exact scripts and cited
source documents; add a manifest with hashes; and use distinct metadata fields
such as \texttt{base\_seed=20260808} and
\texttt{primary\_rng\_seed=20260809}.  Add the clinical-effect factor to the
formal mature-model equations and retain machine-checkable assertions for the
modal totals, no-double-count allocation, public result summaries, and
Tier-C-zero logic.

\subsection*{7. Independent Poisson benefit and harm counts require
justification}

\textbf{Status: likely issue; probably numerically small because the event
probabilities are small.}

\textbf{Location.} Equation~(12), where $S^+$ and $S^-$ are sampled as
independent Poisson counts.

\textbf{Problem.} Beneficial and adverse outcomes arise from the same finite
set of affected incidents and are generally mutually exclusive for an
individual outcome.  Independent Poisson draws are exact only under an
appropriate Poisson-exposure/thinning construction, which is not stated.  With
fixed annual opportunities, a binomial or multinomial model is the direct
construction.

\textbf{Why it matters.} The approximation discards the finite-population and
mutual-exclusion covariance and can inflate predictive variance.  In addition,
the term ``predictive interval'' may be read more broadly than intended: the
calculation adds conditional count noise but not unmodeled annual changes in
case mix, adoption, or system behavior.

\textbf{Suggested fix.} Either justify the Poisson-thinning approximation and
quantify its effect, or simulate pathway-specific affected incidents and
multinomial outcome categories.  Rename the interval ``conditional annual-count
predictive interval'' to delimit what is and is not predicted.

\section*{Editorial, citation, and reporting findings}

\begin{enumerate}
\item \textbf{Headline inconsistency.} The abstract combines the parameter
median $3535$ with the predictive interval $2139$--$5700$, whereas Table~10's
predictive median is $3536$.  Use one internally matched row throughout.  The
quoted result and conclusion already use $3536$.

\item \textbf{Reference \texttt{vertical2007}.} The entry attributes
\emph{The ``Vertical Response Time'': Barriers to Ambulance Response in an
Urban Area} to ``Morrison LJ et al.''  The cited 449-call New York study is by
Silverman RA, Galea S, Blaney S, et al.  Correct the authors and add the journal,
volume, pages, and DOI; the current entry appears to conflate two different
patient-access studies.

\item \textbf{Reference \texttt{pediatric2026}.} Verify and update the journal
title.  The 2026 article record appears under \emph{Circulation: Population
Health and Outcomes}, not \emph{Circulation: Cardiovascular Quality and
Outcomes}.

\item \textbf{List grammar near \texttt{nguyen2024}.} The penultimate evidence
bullet places the word ``and'' after the \texttt{nguyen2024} citation and
immediately before a new list item.  Remove the dangling ``and'' or rewrite the
bullets as complete sentences.

\item \textbf{NEMSIS 2025 anchor.} The cited \texttt{nemsis2025} landing page
does not document the stated $63.636$ million count or the extraction/filter
logic.  Supply a dated data release, query, or archived output sufficient to
verify that number.
\end{enumerate}

\section*{Proofing sweep}

The bundled mechanical scan was run on the compiled mortality manuscript, the
companion economic manuscript, and the mortality source.  Its only mortality
hits were semicolon-capitalization candidates inside compact table cells; spot
checking showed these were false positives.  No duplicated-comma/period,
malformed ``into \ldots{} into'' frame, product-name capitalization, or
\texttt{arctan(x/y)} branch pattern was found.  The high-confidence reference
and list defects identified manually are reported above.

\section*{Checks completed without an additional concrete error}

The principal displayed equations were checked for units and signs.  Equations
for $L_{988}$, $d(p,\beta,t)$, $A_j$, $\lambda_j^{\pm}$, and $L_F$ are
dimensionally consistent as written.  The modal pathway table adds correctly
from its published inputs.  The supplied mortality scripts reproduce all
archived current and mature JSON summaries exactly.  No inverse-trigonometric,
quadrant, coordinate-frame, or branch-convention expression occurs in the
manuscript.  These checks do not resolve the population, endpoint, and overlap
problems above.

\section*{Items needing external verification or unavailable materials}

The bibliographic metadata and quoted results for \texttt{patel2026},
\texttt{nguyen2024}, \texttt{video2025}, and the 2024 CARES summary were
spot-checked against primary records and are broadly consistent, subject to the
population-transport issues already described.  The $63.636$ million 2025
NEMSIS extraction and the author-supplied function/scenario documents remain
unverifiable from the deposited files.  The probabilistic mortality summaries
were independently reproduced from the supplied scripts and JSON input.

\section*{Highest-priority fixes}

\begin{enumerate}
\item Replace ``$L_{\mathrm{ID}}=0$'' with a nonnumeric not-identified statement.
\item Rebuild the current OHCA calculation using outcome-study-compatible
eligible denominators and baseline risks, with time endpoints defined
explicitly.
\item Enforce one mature incident partition across the OHCA, language, and
location/access mechanisms and rerun all mortality and economic results.
\item Track and deposit the scripts and source documents, disambiguate the base
and actual RNG seeds, and write the clinical-effect operation into the formal
model.
\item Relabel Tier-C-zero and add a sensitivity that removes all nonidentified
positive PEC mortality effects.
\item Only after rerunning the model, synchronize the abstract, quoted result,
conclusion, companion economic paper, and archived JSON files.
\end{enumerate}

\end{document}